\chapter{ผลการดำเนินงานด้านตัวชี้วัด}

\section{ตัวชี้วัดของมหาวิทยาลัยที่กลุ่มแผนงานรับมาดำเนินการ}
กลุ่มแผนงานใต้ร่มพระบารมีรับตัวชี้วัดตามยุทธศาสตร์ความเป็นเลิศของมหาวิทยาลัยมาดำเนินการ 7~ตัว
โดยมีตัวชี้วัดที่ 39 เป็นตัวชี้วัดที่ผูกกับภารกิจของกลุ่มแผนงานโดยตรง
ส่วนตัวชี้วัดอื่นเป็นเป้าระดับมหาวิทยาลัยที่กลุ่มแผนงานเป็นผู้สนับสนุนส่วนหนึ่ง

\begin{table}[H]
\centering
\caption{ตัวชี้วัดที่รับมาดำเนินการเทียบเป้าหมายของมหาวิทยาลัย}
\footnotesize
\begin{tabular}{L{1.5cm}L{5.6cm}rrrr}
\toprule
\rowcolor{headerblue}
\textbf{ตัวชี้วัด} & \textbf{รายละเอียด} & \textbf{เป้า มทร.} & \textbf{รับมา} & \textbf{ร้อยละ} & \textbf{โครงการ}\\
\midrule
\rowcolor{rowblue}
\textbf{ที่ 39} & \textbf{องค์ความรู้และเทคโนโลยีที่นำไปใช้ในโครงการหลวงและโครงการตามพระราชดำริ (เป้าของกลุ่มโดยตรง)} & \textbf{60} & \textbf{62} & \textbf{103\%} & \textbf{37}\\
ที่ 35 & คณาจารย์และบุคลากรที่นำเทคโนโลยีไปพัฒนาชุมชนและภาคอุตสาหกรรม & 400 & 288 & 72\% & 46\\
\rowcolor{rowblue}
ที่ 40 & องค์ความรู้และนวัตกรรมที่ยกระดับคุณภาพชีวิตชุมชน & 100 & 45 & 51\% & 33\\
ที่ 36 & แหล่งเรียนรู้ตลอดชีวิตของสังคม & 15 & 4 & 27\% & 4\\
\rowcolor{rowblue}
ที่ 10 & เครือข่ายความร่วมมือกับหน่วยงานภายนอก & 50 & 9 & 18\% & 4\\
ที่ 17 & ทรัพย์สินทางปัญญาที่ยื่นขอจดทะเบียน & 60 & 8 & 13\% & 8\\
\rowcolor{rowblue}
ที่ 38 & สถานประกอบการที่ได้รับการถ่ายทอดเทคโนโลยี & 100 & 8 & 8\% & 7\\
\bottomrule
\end{tabular}
\end{table}

\begin{infobox}
เป้าหมายในคอลัมน์ ``เป้า มทร.'' เป็นเป้าระดับมหาวิทยาลัย กลุ่มแผนงานใต้ร่มพระบารมี
เป็นผู้สนับสนุนส่วนหนึ่งเท่านั้น \textbf{ยกเว้นตัวชี้วัดที่ 39 ซึ่งเป็นเป้าหมายของกลุ่มโดยตรง}
ดังนั้นการที่ตัวชี้วัดที่ 10 17 และ 38 อยู่ในระดับต่ำ จึงไม่ได้แปลว่ามหาวิทยาลัยทำไม่ถึงเป้า
แต่หมายความว่ากลุ่มแผนงานนี้มีส่วนสนับสนุนตัวชี้วัดกลุ่มดังกล่าวน้อย
\end{infobox}

\section{การกระจายตัวชี้วัดตามโครงการหลัก}
เมื่อแยกตัวชี้วัดตามโครงการหลัก พบว่าการแบ่งบทบาทของทั้งสามโครงการมีความชัดเจน

\begin{table}[H]
\centering
\caption{จำนวนตัวชี้วัดที่รับมาดำเนินการ จำแนกตามโครงการหลัก}
\small
\begin{tabular}{L{3.4cm}rrrrrr}
\toprule
\rowcolor{headerblue}
\textbf{โครงการหลัก} & \textbf{ที่ 10} & \textbf{ที่ 17} & \textbf{ที่ 35} & \textbf{ที่ 36} & \textbf{ที่ 38} & \textbf{ที่ 39 / 40}\\
\midrule
ง8-1 ผลักดันเทคโนโลยี & 1 & \textbf{8} & 98 & -- & 2 & 26 / 13\\
\rowcolor{rowblue}
ง8-2 ขับเคลื่อนองค์ความรู้ & 1 & -- & 58 & \textbf{4} & 1 & 13 / 10\\
ง8-3 พัฒนากำลังคน & 7 & -- & \textbf{132} & -- & 5 & 23 / 22\\
\midrule
\rowcolor{headerblue}
\textbf{รวม} & \textbf{9} & \textbf{8} & \textbf{288} & \textbf{4} & \textbf{8} & \textbf{62 / 45}\\
\bottomrule
\end{tabular}
\end{table}

ข้อสังเกตสำคัญสามประการ ประการแรก \textbf{ทรัพย์สินทางปัญญาทั้งหมด} (ตัวชี้วัดที่ 17 จำนวน 8~ผลงาน)
อยู่ในโครงการหลัก ง8-1 ซึ่งเป็นแกนด้านเทคโนโลยี ประการที่สอง \textbf{แหล่งเรียนรู้ทั้งหมด}
(ตัวชี้วัดที่ 36 จำนวน 4~แหล่ง) อยู่ใน ง8-2 ซึ่งเป็นแกนด้านองค์ความรู้ และประการที่สาม
\textbf{การพัฒนากำลังคน} (ตัวชี้วัดที่ 35 จำนวน 132~คน) กระจุกตัวใน ง8-3 ตามชื่อโครงการ
สะท้อนว่าการออกแบบโครงสร้าง ง8 สอดคล้องกับผลที่แต่ละโครงการรับผิดชอบจริง

\section{สถานะการอนุมัติโครงการ}
ณ วันที่จัดทำรายงาน โครงการที่ได้รับอนุมัติและโอนงบประมาณแล้วมีจำนวน 59~โครงการ
จากทั้งหมด 61~โครงการ คิดเป็นร้อยละ~97 ส่วนที่เหลืออีก 2 โครงการอยู่ระหว่างการปรับแก้ไขรายละเอียดโครงการ
